
\documentclass[11pt,a4paper]{article}
\usepackage[utf8]{inputenc}
\usepackage[T1]{fontenc}
\usepackage[french]{babel}
\usepackage{lmodern}
\usepackage{microtype}
\usepackage[a4paper,margin=1.65cm,headheight=15pt]{geometry}
\usepackage{amsmath,amssymb,mathtools}
\usepackage{array,tabularx,multicol,booktabs}
\usepackage{enumitem}
\usepackage[most]{tcolorbox}
\usepackage{xcolor}
\usepackage{tikz}
\usetikzlibrary{arrows.meta,calc,positioning,angles,quotes}
\usepackage{pdflscape}
\usepackage{fancyhdr}
\usepackage{hyperref}
\usepackage{needspace}
\usepackage{pagecolor}

\definecolor{fondpage}{RGB}{247,248,250}
\definecolor{encre}{RGB}{31,35,40}
\definecolor{bleu}{RGB}{40,91,135}
\definecolor{vert}{RGB}{55,125,92}
\definecolor{orange}{RGB}{178,105,42}
\definecolor{violet}{RGB}{101,77,145}
\definecolor{rouge}{RGB}{170,72,72}
\definecolor{gris}{RGB}{86,91,99}
\definecolor{bleufonce}{RGB}{28,55,120}
\definecolor{bleuclair}{RGB}{238,243,255}
\definecolor{grisclair}{RGB}{245,245,245}
\definecolor{tableline}{RGB}{45,45,45}
\definecolor{softgray}{RGB}{246,247,249}
\definecolor{deepblue}{RGB}{25,62,115}
\definecolor{accent}{RGB}{20,82,130}

\hypersetup{colorlinks=true,linkcolor=bleufonce,urlcolor=bleufonce,
pdftitle={Parcours de revision -- Produit scalaire -- Premiere specialite},pdfauthor={}}

\setlength{\parindent}{0pt}
\setlength{\parskip}{0.25em}
\renewcommand{\arraystretch}{1.13}
\setlength{\tabcolsep}{4pt}
\setlist[itemize]{leftmargin=1.25em,itemsep=0.15em,topsep=0.2em}
\setlist[enumerate,1]{label=\textbf{\arabic*.}, leftmargin=1.05cm, itemsep=0.35em, topsep=0.2em}
\setlist[enumerate,2]{label=\textbf{\alph*.}, leftmargin=0.85cm, itemsep=0.25em, topsep=0.15em}

\pagestyle{fancy}
\fancyhf{}
\cfoot{\thepage}

\newcommand{\R}{\mathbb{R}}
\newcommand{\N}{\mathbb{N}}
\newcommand{\vect}[1]{\overrightarrow{#1}}
\newcommand{\norme}[1]{\left\lVert #1\right\rVert}
\newcommand{\sourceq}[1]{ {\scriptsize\textcolor{bleufonce}{(site Premiere q#1)}}}

\newcounter{question}
\newcounter{reponse}

\newenvironment{blocquestions}[1]{%
  \begin{tcolorbox}[enhanced,breakable,colback=bleuclair,colframe=bleufonce,boxrule=0.6pt,arc=2mm,left=2mm,right=2mm,top=2mm,bottom=1.5mm,title={\bfseries #1},coltitle=white,colbacktitle=bleufonce,before skip=3mm,after skip=3mm, fontupper=\large]
  \large
  \begin{enumerate}[leftmargin=*,label=\textbf{\arabic*.},itemsep=0.82em,topsep=0.5em]
  \setcounter{enumi}{\value{question}}
}{%
  \setcounter{question}{\value{enumi}}
  \end{enumerate}
  \end{tcolorbox}
}

\newenvironment{blocreponses}[1]{%
  \begin{tcolorbox}[enhanced,breakable,colback=bleuclair,colframe=bleufonce,boxrule=0.6pt,arc=2mm,left=2mm,right=2mm,top=2mm,bottom=1.5mm,title={\bfseries #1},coltitle=white,colbacktitle=bleufonce,before skip=3mm,after skip=3mm, fontupper=\large]
  \large
  \begin{enumerate}[leftmargin=*,label=\textbf{\arabic*.},itemsep=0.42em,topsep=0.5em]
  \setcounter{enumi}{\value{reponse}}
}{%
  \setcounter{reponse}{\value{enumi}}
  \end{enumerate}
  \end{tcolorbox}
}

\newtcolorbox{correctionbox}[1]{enhanced,breakable,colback=grisclair,colframe=black!55,boxrule=0.55pt,arc=2mm,left=2mm,right=2mm,top=2mm,bottom=1.5mm,title=\textbf{#1},coltitle=black,colbacktitle=black!12,before skip=2mm,after skip=2mm}
\newtcolorbox{methodebox}[1]{enhanced,breakable,colback=white,colframe=bleu!70!black,boxrule=0.55pt,arc=2mm,left=2mm,right=2mm,top=2mm,bottom=1.5mm,title=\textbf{#1},coltitle=white,colbacktitle=bleu!70!black,before skip=2mm,after skip=2mm}

\newcommand{\titreSujet}[2]{%
  \subsection*{#1}
  \addcontentsline{toc}{subsection}{#1}
  \textit{Référence / inspiration : #2}\par\medskip\hrule\medskip
}
\newcommand{\qcm}[4]{%
\begin{center}\begin{tabularx}{0.96\linewidth}{@{}XX@{}}
\textbf{a.} #1 & \textbf{b.} #2\\[1mm]
\textbf{c.} #3 & \textbf{d.} #4
\end{tabularx}\end{center}}

\tcbset{enhanced,arc=2mm,outer arc=2mm,boxrule=0.42pt,left=1.15mm,right=1.15mm,top=0.75mm,bottom=0.75mm,boxsep=0.35mm,before skip=0.8mm,after skip=0.8mm,fonttitle=\bfseries\footnotesize,coltitle=encre,before upper=\raggedright\fontsize{7.05}{7.55}\selectfont}
\newtcolorbox{fichebox}[3][]{colback=white,colframe=#2!72!black,colbacktitle=#2!18,coltitle=#2!50!black,borderline west={0.9mm}{0pt}{#2!78!black},title={#3},drop fuzzy shadow=black!5,#1}
\newcommand{\tagcol}[2]{\begin{tcolorbox}[enhanced,colback=#1!11,colframe=#1!45!black,boxrule=0.42pt,arc=2mm,left=1mm,right=1mm,top=0.45mm,bottom=0.45mm,before skip=0mm,after skip=0.85mm,fontupper=\bfseries\scriptsize,colupper=#1!55!black]\centering #2\end{tcolorbox}}
\newtcolorbox{panel}[1]{colback=fondpage,colframe=fondpage,boxrule=0pt,arc=0mm,left=0mm,right=0mm,top=0mm,bottom=0mm,height=168mm,valign=top,before skip=0mm,after skip=0mm}

% Mini-schémas TikZ inspirés des fiches visuelles de Première
\newcommand{\schemaAngleProduit}{%
\begin{center}
\begin{tikzpicture}[scale=0.48,>=Latex,line cap=round]
  \coordinate (A) at (0,0);
  \coordinate (B) at (2.45,0);
  \coordinate (C) at (1.55,1.35);
  \draw[->,very thick,bleu!85!black] (A)--(B) node[midway,below=2pt] {\scriptsize $\vect{AB}$};
  \draw[->,very thick,orange!88!black] (A)--(C) node[midway,above left=1pt] {\scriptsize $\vect{AC}$};
  \pic[draw=gris,thick,angle radius=0.46cm,"{\scriptsize $\theta$}",angle eccentricity=1.45] {angle=B--A--C};
  \fill (A) circle (1.2pt) node[below left=1pt] {\scriptsize $A$};
  \fill (B) circle (1.2pt) node[below right=1pt] {\scriptsize $B$};
  \fill (C) circle (1.2pt) node[above=1pt] {\scriptsize $C$};
\end{tikzpicture}
\end{center}}
\newcommand{\schemaProjectionProduit}{%
\begin{center}
\begin{tikzpicture}[scale=0.50,>=Latex,line cap=round]
  \coordinate (A) at (0,0);
  \coordinate (B) at (3.65,0);
  \coordinate (H) at (2.35,0);
  \coordinate (C) at (2.35,1.55);

  % base et vecteurs
  \draw[->,very thick,bleu!85!black] (A)--(B) node[midway,below=6pt] {\scriptsize $\vect{AB}$};
  \draw[->,very thick,orange!88!black] (A)--(C) node[midway,above left=1pt] {\scriptsize $\vect{AC}$};

  % projection
  \draw[dashed,gris,thick] (C)--(H);
  \draw[gris,thick] ($(H)+(-0.18,0)$)--++(0,0.18)--++(0.18,0);

  % points
  \fill (A) circle (1.2pt) node[below left=1pt] {\scriptsize $A$};
  \fill (B) circle (1.2pt) node[below right=1pt] {\scriptsize $B$};
  \fill (C) circle (1.2pt) node[above=1pt] {\scriptsize $C$};
  \fill (H) circle (1.2pt) node[below=2pt] {\scriptsize $H$};
\end{tikzpicture}
\end{center}}
\newcommand{\schemaNormalDroite}{%
\begin{center}
\begin{tikzpicture}[scale=0.42,>=Latex,line cap=round]
  % Droite : segment simple, sans flèche.
  \draw[thick,bleu!80!black,line cap=butt] (-1.25,-.48)--(2.25,.86);
  \node[bleu!80!black,right=2pt] at (2.25,.86) {\scriptsize $d$};
  \coordinate (P) at (.5,.25);
  \draw[->,thick,rouge!85!black] (P)--++(-.55,1.45) node[above] {\scriptsize $\vec n$};
\end{tikzpicture}
\end{center}}
\newcommand{\schemaCercleDiametre}{%
\begin{center}
\begin{tikzpicture}[scale=0.48,line cap=round]
  \coordinate (A) at (-1.60,0);
  \coordinate (B) at (1.60,0);
  \coordinate (O) at (0,0);
  \coordinate (M) at (0.42,1.544);

  % Cercle de diamètre [AB]
  \draw[thick,bleu!80!black] (O) circle (1.60);

  % Diamètre et triangle inscrit
  \draw[thick] (A)--(B);
  \draw[thick] (A)--(M)--(B);

  % Angle droit en M
  \pic[draw=rouge!70!black,thick,angle radius=.28cm] {right angle=A--M--B};

  % Points et labels
  \fill (A) circle (1.15pt) node[below left=2pt] {\scriptsize $A$};
  \fill (B) circle (1.15pt) node[below right=2pt] {\scriptsize $B$};
  \fill (M) circle (1.15pt) node[above=2pt] {\scriptsize $M$};
\end{tikzpicture}
\end{center}}
\newcommand{\schemaTangenteCercle}{%
\begin{center}
\begin{tikzpicture}[scale=0.42,>=Latex,line cap=round]
  \coordinate (I) at (0,0);\coordinate (T) at (1.45,.55);
  \draw[thick,bleu!80!black] (I) circle (1.55);
  \draw[thick,rouge!85!black] ($(T)+(-.55,1.45)$)--($(T)+(.55,-1.45)$) node[below] {\scriptsize tangente};
  \draw[->,thick,orange!90!black] (I)--(T) node[midway,below] {\scriptsize rayon};
  \fill (I) circle (1.2pt) node[left] {\scriptsize $I$};\fill (T) circle (1.2pt) node[right] {\scriptsize $T$};
  \draw[gris] (T) ++(-.07,.17)--++(.17,.07)--++(.07,-.17);
\end{tikzpicture}
\end{center}}

\newcommand{\titrepage}{\begin{tcolorbox}[colback=white,colframe=bleu!70!black,boxrule=0.65pt,arc=3mm,left=2.5mm,right=2.5mm,top=1.15mm,bottom=1.15mm,before skip=0mm,after skip=1.4mm,drop fuzzy shadow=black!10]
{\Large\bfseries Parcours de révision -- Produit scalaire}\hfill {\normalsize Première spécialité -- sans calculatrice}
\end{tcolorbox}}

\newcommand{\ficheRecap}{%
\begin{landscape}
\pagecolor{fondpage}
\thispagestyle{empty}
\titrepage
\begin{tabularx}{\linewidth}{@{}X@{\hspace{1.6mm}}X@{\hspace{1.6mm}}X@{\hspace{1.6mm}}X@{}}

\begin{panel}{}
\tagcol{gris}{Rappels de Seconde}
\begin{fichebox}{gris}{Coordonnées de $\vect{AB}$}
Si $A(x_A;y_A)$ et $B(x_B;y_B)$, alors
\[
\vect{AB}=\begin{pmatrix}x_B-x_A\\ y_B-y_A\end{pmatrix}.
\]
\end{fichebox}
\begin{fichebox}{gris}{Distance}
Dans un repère orthonormé :
\[
AB=\sqrt{(x_B-x_A)^2+(y_B-y_A)^2}.
\]
Donc $AB^2=(x_B-x_A)^2+(y_B-y_A)^2$.
\end{fichebox}
\begin{fichebox}{gris}{Milieu d'un segment}
Le milieu $I$ de $[AB]$ a pour coordonnées :
\[
I\left(\frac{x_A+x_B}{2};\frac{y_A+y_B}{2}\right).
\]
\end{fichebox}
\begin{fichebox}{gris}{À utiliser avant le produit scalaire}
Pour calculer $\vect{AB}\cdot\vect{AC}$, on commence souvent par calculer les coordonnées de $\vect{AB}$ et de $\vect{AC}$.
\end{fichebox}
\end{panel}
&
\begin{panel}{}
\tagcol{bleu}{Définitions et calculs}
\begin{fichebox}{bleu}{Produit scalaire en coordonnées}
Dans un repère orthonormé, si
\[
\vec u=\begin{pmatrix}x\\y\end{pmatrix},\qquad
\vec v=\begin{pmatrix}x'\\y'\end{pmatrix},
\]
alors
\[
\vec u\cdot\vec v=xx'+yy'.
\]
\end{fichebox}
\begin{fichebox}{bleu}{Avec les longueurs}
\schemaAngleProduit
\[
\vect{AB}\cdot\vect{AC}=AB\times AC\times\cos(\widehat{BAC}).
\]
En particulier : \(\vec u\cdot\vec u=\norme{\vec u}^{2}\).
\end{fichebox}
\begin{fichebox}{bleu}{Propriétés de calcul}
Symétrie : \(\vec u\cdot\vec v=\vec v\cdot\vec u\).

Bilinéarité :
\[
(\vec u+\vec v)\cdot\vec w=\vec u\cdot\vec w+\vec v\cdot\vec w,
\]
\[
(k\vec u)\cdot\vec v=k(\vec u\cdot\vec v).
\]
Identité utile :
\[
\norme{\vec u+\vec v}^{2}=\norme{\vec u}^{2}+2\vec u\cdot\vec v+\norme{\vec v}^{2}.
\]
\end{fichebox}
\begin{fichebox}{bleu}{Al-Kashi}
Dans un triangle \(ABC\) :
\[
BC^2=AB^2+AC^2-2\,AB\times AC\cos(\widehat{BAC}).
\]
\end{fichebox}
\end{panel}
&
\begin{panel}{}
\tagcol{vert}{Orthogonalité et angle}
\begin{fichebox}{vert}{Orthogonalité}
Deux vecteurs non nuls \(\vec u\) et \(\vec v\) sont orthogonaux si et seulement si :
\[
\vec u\cdot\vec v=0.
\]
Donc deux droites sont perpendiculaires si leurs vecteurs directeurs ont un produit scalaire nul.
\end{fichebox}
\begin{fichebox}{vert}{Calculer un angle}
Si \(\vec u\) et \(\vec v\) sont non nuls :
\[
\cos\theta=\frac{\vec u\cdot\vec v}{\norme{\vec u}\,\norme{\vec v}}.
\]
Le signe du produit scalaire renseigne sur l’angle : aigu si positif, droit si nul, obtus si négatif.
\end{fichebox}
\begin{fichebox}{vert}{Projeté orthogonal}
\schemaProjectionProduit
Si \(H\) est le projeté orthogonal de \(C\) sur \((AB)\), alors :
\[
\vect{AB}\cdot\vect{AC}=\vect{AB}\cdot\vect{AH}.
\]
\end{fichebox}
\end{panel}
&
\begin{panel}{}
\tagcol{orange}{Applications repérées}
\begin{fichebox}{orange}{Vecteur normal}
\schemaNormalDroite
Pour une droite \(ax+by+c=0\), un vecteur normal est
\[
\vec n=\begin{pmatrix}a\\b\end{pmatrix}.
\]
Un vecteur directeur possible est \(\vec d=\begin{pmatrix}-b\\a\end{pmatrix}\).
\end{fichebox}
\begin{fichebox}{orange}{Équation de cercle}
Le cercle de centre \(I(x_I;y_I)\) et de rayon \(r\) a pour équation :
\[
(x-x_I)^2+(y-y_I)^2=r^2.
\]
\end{fichebox}
\begin{fichebox}{orange}{Cercle de diamètre}
\schemaCercleDiametre
Un point \(M\) appartient au cercle de diamètre \([AB]\) si et seulement si :
\[
\vect{MA}\cdot\vect{MB}=0.
\]
\end{fichebox}
\end{panel}
\end{tabularx}
\clearpage
\nopagecolor
\end{landscape}}

\begin{document}
\begin{center}
{\LARGE\bfseries Parcours de révision -- Produit scalaire}\\[1mm]
{\large Première générale -- spécialité mathématiques}\\[1mm]
{\normalsize Sans calculatrice}
\end{center}
\tableofcontents\newpage
\phantomsection\addcontentsline{toc}{section}{Fiche récapitulative}
\ficheRecap
\section{Questions flash}
\setcounter{question}{0}
\begin{blocquestions}{Questions flash 1 à 12}
\item Dans un repère orthonormé, calculer le produit scalaire de $\vec u(2;-1)$ et $\vec v(3;4)$.
\item Les vecteurs $\vec u(1;2)$ et $\vec v(4;-2)$ sont-ils orthogonaux ?
\item Calculer $\overrightarrow{AB}$ pour $A(1;3)$ et $B(5;-1)$.
\item Calculer $\overrightarrow{AB}\cdot\overrightarrow{AC}$ pour $A(0;0)$, $B(3;1)$ et $C(2;-6)$.
\item Si $\vec u\cdot\vec v=0$ et $\vec u,\vec v$ sont non nuls, que peut-on dire de $\vec u$ et $\vec v$ ?
\item Calculer $\norme{\vec u}$ pour $\vec u(6;8)$.
\item Si $\vec u\cdot\vec v=6$, $\norme{\vec u}=2$ et $\norme{\vec v}=3$, calculer le cosinus de l'angle formé par $\vec u$ et $\vec v$.
\item Un angle a un cosinus égal à $0$. Que peut-on dire de cet angle ?
\item Donner un vecteur normal à la droite $3x-2y+5=0$.
\item Donner un vecteur directeur de la droite $3x-2y+5=0$.
\item Le point $A(1;2)$ appartient-il à la droite $2x+y-4=0$ ?
\item Donner l’équation du cercle de centre $I(2;-1)$ et de rayon $3$.
\end{blocquestions}
\begin{blocquestions}{Questions flash 13 à 24}
\item Le point $M(5;-1)$ appartient-il au cercle $(x-2)^2+(y+1)^2=9$ ?
\item Déterminer le centre et le rayon du cercle $(x+1)^2+(y-4)^2=25$.
\item Si $M$ appartient au cercle de diamètre $[AB]$, que vaut $\overrightarrow{MA}\cdot\overrightarrow{MB}$ ?
\item Pour un cercle de centre $I$, quel vecteur est normal à la tangente en $T$ ?
\item Calculer $\vec u\cdot\vec u$ si $\norme{\vec u}=5$.
\item Si $\vec u(2;m)$ est orthogonal à $\vec v(3;6)$, déterminer $m$.
\item Les droites de vecteurs directeurs $\vec u(1;4)$ et $\vec v(8;-2)$ sont-elles perpendiculaires ?
\item Calculer $AB$ pour $A(-1;2)$ et $B(3;5)$.
\item Calculer $\overrightarrow{AB}\cdot\overrightarrow{AC}$ pour $A(1;1)$, $B(4;1)$, $C(1;5)$.
\item Que signifie un produit scalaire strictement positif entre deux vecteurs non nuls ?
\item Que signifie un produit scalaire strictement négatif entre deux vecteurs non nuls ?
\item Donner un vecteur normal à la droite $x+4y-7=0$.
\end{blocquestions}
\begin{blocquestions}{Questions flash 25 à 36}
\item Écrire une équation de la droite de vecteur normal $\vec n(2;5)$ passant par $A(1;-1)$.
\item Un vecteur normal à une droite est $\vec n(1;-3)$. Donner un vecteur directeur possible.
\item Déterminer si $A(0;0)$, $B(4;2)$ et $C(1;-2)$ forment un angle droit en $A$.
\item Donner l’équation du cercle de diamètre $[AB]$ avec $A(0;0)$ et $B(4;0)$.
\item Le cercle $(x-1)^2+(y+2)^2=10$ a-t-il pour rayon $\sqrt{10}$ ?
\item Si $\vec u\cdot\vec v=\norme{\vec u}\norme{\vec v}$, que vaut l’angle entre les vecteurs ?
\item Si $\vec u\cdot\vec v=-\norme{\vec u}\norme{\vec v}$, que vaut l’angle entre les vecteurs ?
\item Calculer $\overrightarrow{BA}\cdot\overrightarrow{BC}$ pour $A(1;0)$, $B(0;0)$, $C(0;2)$.
\item Donner une condition vectorielle pour que $M$ soit sur le cercle de diamètre $[AB]$.
\item Dans un triangle, comment peut-on montrer qu’un angle est droit avec le produit scalaire ?
\item Pour la droite $2x-y+1=0$, le vecteur $(2;-1)$ est-il normal ou directeur ?
\item Calculer $\vec u\cdot\vec v$ si $\norme{\vec u}=4$, $\norme{\vec v}=5$ et l’angle vaut $60^\circ$.
\end{blocquestions}

\section{QCM d’automatismes}
Pour chaque question, une seule réponse est correcte.
\begin{enumerate}
\item Si $\vec u(2;3)$ et $\vec v(-1;4)$, alors $\vec u\cdot\vec v=$
\qcm{$10$}{$-2$}{$14$}{$0$}
\item Les vecteurs $\vec u(5;1)$ et $\vec v(1;-5)$ sont :
\qcm{colinéaires}{orthogonaux}{égaux}{opposés}
\item Une droite d’équation $-2x+3y+1=0$ a pour vecteur normal :
\qcm{$(3;2)$}{$(-2;3)$}{$(2;3)$}{$(3;-2)$}
\item Le cercle de centre $I(1;2)$ et de rayon $4$ a pour équation :
\qcm{$(x+1)^2+(y+2)^2=16$}{$(x-1)^2+(y-2)^2=16$}{$(x-1)^2+(y-2)^2=4$}{$(x+1)^2+(y+2)^2=4$}
\item Si $M$ appartient au cercle de diamètre $[AB]$, alors :
\qcm{$\vect{MA}\cdot\vect{MB}=0$}{$\vect{MA}\cdot\vect{MB}=1$}{$MA=MB$}{$M$ est le milieu de $[AB]$}
\item Si $\vec u\cdot\vec v<0$, alors l’angle entre $\vec u$ et $\vec v$ est :
\qcm{aigu}{droit}{obtus}{nul}
\item La tangente à un cercle en $T$ est perpendiculaire :
\qcm{au diamètre horizontal}{au rayon $[IT]$}{à l’axe des ordonnées}{à toute corde}
\item Pour $A(1;2)$ et $B(4;6)$, la longueur $AB$ vaut :
\qcm{$5$}{$7$}{$\sqrt{5}$}{$25$}
\end{enumerate}

\section{Sujets type bac / E3C}

\titreSujet{Sujet 1 -- Triangle, angle et orthogonalité}{inspiré de sujets E3C de géométrie repérée et produit scalaire}
Dans un repère orthonormé, on considère les points :
\[
A(1;2),\qquad B(5;4),\qquad C(3;-2).
\]
\begin{center}
\begin{tikzpicture}[scale=0.58,>=Latex]
\draw[step=1,black!10] (-1,-3) grid (7,5);
\draw[->] (-1.2,0)--(7.2,0) node[right] {$x$};
\draw[->] (0,-3.2)--(0,5.2) node[above] {$y$};
\coordinate (A) at (1,2); \coordinate (B) at (5,4); \coordinate (C) at (3,-2);
\draw[thick] (A)--(B)--(C)--cycle;
\fill (A) circle (2pt) node[above left] {$A$};\fill (B) circle (2pt) node[above right] {$B$};\fill (C) circle (2pt) node[below right] {$C$};
\end{tikzpicture}
\end{center}
\begin{enumerate}
\item Calculer les coordonnées des vecteurs $\overrightarrow{AB}$ et $\overrightarrow{AC}$.
\item Calculer le produit scalaire $\overrightarrow{AB}\cdot\overrightarrow{AC}$.
\item Les droites $(AB)$ et $(AC)$ sont-elles perpendiculaires ?
\item Calculer les longueurs $AB$ et $AC$.
\item En déduire la valeur exacte de $\cos(\widehat{BAC})$.
\item Le triangle $ABC$ est-il rectangle ? Justifier.
\item Calculer $\overrightarrow{BA}\cdot\overrightarrow{BC}$ et interpréter le signe obtenu.
\end{enumerate}

\titreSujet{Sujet 2 -- Droite, vecteur normal et projeté orthogonal}{inspiré d’exercices E3C sur vecteur normal et équation cartésienne}
Dans un repère orthonormé, on considère la droite $d$ d’équation :
\[
2x-y-3=0
\]
et le point $A(4;1)$.
\begin{center}
\begin{tikzpicture}[scale=0.7,>=Latex]
\draw[step=1,black!10] (-1,-1) grid (6,5);
\draw[->] (-1.2,0)--(6.2,0) node[right] {$x$};
\draw[->] (0,-1.2)--(0,5.2) node[above] {$y$};
\draw[domain=1:4,thick,bleu!70!black] plot (\x,{2*\x-3}) node[right] {$d$};
\fill (4,1) circle (2pt) node[below right] {$A$};
\draw[->,rouge!70!black,thick] (4,1)--(5.3,0.35) node[right] {$\vec n$};
\end{tikzpicture}
\end{center}
\begin{enumerate}
\item Donner un vecteur normal à la droite $d$.
\item Justifier que la droite $\Delta$, perpendiculaire à $d$ et passant par $A$, a pour vecteur directeur $\vec n(2;-1)$.
\item Déterminer une équation paramétrique de $\Delta$.
\item Déterminer les coordonnées du point $H$, intersection de $d$ et de $\Delta$.
\item Justifier que $H$ est le projeté orthogonal de $A$ sur $d$.
\item Calculer la distance $AH$.
\item Déterminer une équation cartésienne de la droite parallèle à $d$ passant par $A$.
\end{enumerate}

\titreSujet{Sujet 3 -- Cercle, diamètre et tangente}{inspiré d’exercices E3C sur équation de cercle, diamètre et tangente}
Dans un repère orthonormé, on considère les points :
\[
A(-1;2),\qquad B(5;4).
\]
On note $\mathcal C$ le cercle de diamètre $[AB]$.
\begin{center}
\begin{tikzpicture}[scale=0.58,>=Latex]
\draw[step=1,black!10] (-3,-1) grid (7,7);
\draw[->] (-3.2,0)--(7.2,0) node[right] {$x$};
\draw[->] (0,-1.2)--(0,7.2) node[above] {$y$};
\coordinate (A) at (-1,2); \coordinate (B) at (5,4); \coordinate (I) at (2,3); \coordinate (T) at (1,0);
\draw[thick] (I) circle ({sqrt(10)});
\draw[thick] (A)--(B);
\draw[rouge!70!black,thick] (-2.0,1)--(4.6,-1.2) node[right] {tangente en $T$};
\fill (A) circle (2pt) node[above left] {$A$};\fill (B) circle (2pt) node[above right] {$B$};\fill (I) circle (2pt) node[above] {$I$};\fill (T) circle (2pt) node[below] {$T$};
\draw[dashed] (I)--(T);
\end{tikzpicture}
\end{center}
\begin{enumerate}
\item Déterminer les coordonnées du centre $I$ du cercle $\mathcal C$.
\item Calculer le rayon de $\mathcal C$.
\item Déterminer une équation cartésienne de $\mathcal C$.
\item Soit $M(x;y)$. Montrer que $M\in \mathcal C$ équivaut à :
\[
\overrightarrow{MA}\cdot\overrightarrow{MB}=0.
\]
\item Vérifier que le point $T(1;0)$ appartient au cercle $\mathcal C$.
\item Déterminer un vecteur normal à la tangente au cercle en $T$.
\item En déduire une équation de la tangente au cercle $\mathcal C$ en $T$.
\end{enumerate}
\end{document}
